% !TEX root = ../../proposal.tex

\section{Phase 2: Explaining the Composability Gap}
\label{sec:phase_2}

Phase~1 characterised the composability gap across concept regimes, and Phase~2 investigates its cause.

Compositional diffusion methods combine pretrained models at inference time by modifying the score field rather than retraining a joint model.
\citet{du2024compositional} derive this approach from a probabilistic perspective, showing that under conditional independence assumptions, the score of a joint distribution can be expressed as a combination of conditional scores relative to an unconditional model.
In practice, they further show that the behaviour of this composed system depends on inference-time choices such as guidance weighting and sampling dynamics, which can improve stability and sample quality but do not alter the underlying assumptions required for successful composition.

Subsequent theoretical work clarifies when such composition should succeed.
\citet{mechanisms2025projective} show that successful composition depends on the existence of a feature space in which the component concepts act on separable parts of the representation.
In that space, the relevant conditionals become approximately factorised, potentially only after an unknown transformation $A$ of the data.

Instead of assuming that a suitably factorised feature space already exists, Phase~2 asks whether
learning such a representation improves compositional generation under a shared downstream
conditional latent diffusion setup.
A \emph{representation problem} occurs when the model lacks a feature space in which the target concepts are separated enough to support factorised conditionals.
Phase~2 focuses specifically on whether better factorisation by itself improves composition, rather than on ruling out all other possible inference-time improvements.
This leads to a central claim:
\begin{center}
{\setlength{\fboxsep}{6pt}%
\fbox{%
\parbox{\dimexpr 0.96\linewidth - 2\fboxsep - 2\fboxrule\relax}{%
\textbf{{If representation quality is a limiting factor for composition, then more factorised latent
spaces should yield more reliable compositional generation under the same downstream conditional
diffusion architecture.}
}
}}}
\end{center}

H2 formalises this claim.
If the representation is improved, composition success should increase when the same conditional
latent diffusion architecture, conditioning scheme, and compositional inference procedure are used.

This distinction is difficult to test directly in a large pretrained foundation model such as Stable Diffusion \citep{rombach2022high}, where the learned representation and the inference-time composition rule are tightly coupled and cannot be varied independently without substantial confounds.
To test H2 cleanly, Phase~2 moves to a controlled setting in which representation quality can be
varied while the composition rule is kept fixed.

To test this hypothesis, we investigate VAE-based methods for learning a factorised latent space that
approximates the transformation required for factorised conditionals, and then evaluate each learned
latent space with the same downstream conditional latent diffusion architecture.


\subsection{Dataset Selection}
\label{subsec:p2_testbeds}

To test the claim that compositional failure is primarily a representation problem, Phase~2
requires a domain in which representation quality can be manipulated independently of the
inference rule. We therefore use the standard dSprites dataset~\citep{dsprites17} and a
correlated variant~\citep{roth2023disentanglement}. Within the correlated setting, we then
introduce a representation intervention to test whether better factorisation restores
composition.

The dSprites dataset is well suited for this phase because it is simple, controlled, and
fully interpretable. Each image contains a single object whose appearance varies along a
small number of known attributes: shape, scale, orientation, and x/y position. 
In other words, the task defined by this dataset is not composition across distinct objects in a scene, but recombination of multiple attributes of a single object.
If held-out recombination improves here, that improvement can be likened more to improved quality of the learned representation.
Because the true factors of variation are known, dSprites also allows
held-out composition analysis and disentanglement metrics \citep{disentengaled_representation_learning, on_disentangled_representation_learning} to be evaluated together.

% \begin{table}[H]
%     \centering
%     \small
%     \caption{Factor structure of dSprites in Phase~2. Color is fixed; Regime~C2 later supervises
%     the scale and orientation factors directly.}
%     \label{tab:dsprites_factor_table}
%     \begin{tabularx}{0.82\linewidth}{@{}l c X@{}}
%         \toprule
%         Factor & \# values & Notes \\
%         \midrule
%         color & 1 & fixed white foreground \\
%         shape & 3 & square, ellipse, heart \\
%         scale & 6 & physical values 0.5 to 1.0 (e.g classes 0-5 correspond to physical values 0.5, 0.6, 0.7, 0.8, 0.9)\\
%         orientation & 40 & angles from 0 to $2\pi$ radians in uniform $\approx 9.2^\circ$ steps \\
%         posX & 32 & horizontal position \\
%         posY & 32 & vertical position \\
%         \midrule
%         \multicolumn{3}{@{}l@{}}{\textbf{Total images:} $3 \times 6 \times 40 \times 32 \times 32 = 737{,}280$} \\
%         \bottomrule
%     \end{tabularx}
% \end{table}

\begin{figure}[H]
    \centering
    \includegraphics[width=0.26\textwidth]{chapters/research_method/media/dsprites_standard.png}
    \caption{Illustrative samples from the standard dSprites dataset~\citep{dsprites17}. The
    dataset varies a single sprite along known attributes including shape, scale,
    orientation, and position, making it a clean controlled setting for testing attribute
    recombination.}
    \label{fig:p2_dsprites_standard}
\end{figure}

We use standard dSprites as the Phase~2 positive control due to its independently sampled attributes. 
Formally, this is captured by $p(\mathbf{z}) = \prod_i p(z_i)$, meaning factors like orientation and scale do not constrain one another.
To test for held-out recombination, we remove a single orientation-scale combination from the training set.
This setup determines if a model can recover a withheld combination in the absence of shortcut correlations, as the training distribution itself encourages no such dependencies.
If a model fails to recover a withheld orientation-scale combination, that failure reflects the learned representation itself.

\paragraph{Correlated dSprites:}
Correlated dSprites preserves the same visual domain as standard dSprites while modifying the sampling procedure.
We intentionally make orientation and scale correlated in part of the training set.
For a selected subset of the training data, we sample scale from an orientation-dependent distribution rather than independently of orientation. 
Specifically, once an orientation is fixed, one scale is given higher probability than the others. 
This makes some orientation-scale combinations appear much more often than others, while some are rare or missing entirely. 
Following \citet{traeuble2021disentangled}, this yields a controlled setting in which orientation and scale are no longer sampled independently.
Figure~\ref{fig:p2_dsprites_structured} illustrates this visually: once some combinations appear much more often than others, recombination is no longer learned from an evenly represented set of factor combinations.
\begin{figure}[H]
    \centering
    \includegraphics[width=\textwidth]{chapters/research_method/media/dsprites_structured_exclusions.png}
    \caption{Illustrative dSprites train/test exclusions and their effect on latent usage in dSprites \citep{montero2021role}. \textbf{(a)} Example image pairs for three
    generalisation settings: recombination to element, recombination to range, and
    extrapolation. Within each setting, the two columns show representative \emph{training}
    and \emph{test} samples, while the rows indicate the \emph{reference} image, the
    \emph{transform} to be applied, and the resulting \emph{transformed} image. The labels
    beneath the images indicate which factor is being changed, such as shape, orientation, or
    x-position. \textbf{(b)} Corresponding latent-factor association patterns for the same
    three settings. Here, rows correspond to the ground-truth factors (shape, scale,
    orientation, posX, posY), columns correspond to latent dimensions, and brighter cells
    indicate stronger usage or alignment of a latent dimension with a factor. Taken together,
    the figure illustrates how structured exclusions create uneven support over factor
    combinations and alter how factors are represented, making recombination and
    generalisation harder than in the standard factorised dSprites setting.}
    \label{fig:p2_dsprites_structured}
\end{figure}

In standard dSprites, a left-oriented sprite can appear at both small and large scales, and the same is true for a right-oriented sprite.
In constrast, the correlated dSprites training set may contain many examples of left+small and right+large, while left+large and right+small are rare or absent. 
The factors are therefore coupled through sampling rather than varying independently. 
This matters because the model can learn the shortcut that one attribute predicts the other, instead of learning orientation and scale as separate attributes. 
Correlated dSprites keeps the visual domain unchanged while altering the training distribution, providing a controlled setting to study representation learning under correlated attributes.

\subsection{Experiment Design}
\label{subsec:p2_models}

Our experimental design consists of four regimes that progressively test the representation hypothesis under controlled changes to the training distribution and model class. 
Regimes~A and~B establish the basic comparison by keeping the VAE fixed and changing only the training support, from factorised to correlated.
Based on prior work on disentanglement under correlated training data \citep{traeuble2021disentangled, funke2022disentanglement,roth2023disentanglement}, we expect Regime~B to fail on correlated dSprites. 
Regime~C then keeps the correlated setting fixed and asks whether stronger disentanglement objectives improve held-out recombination, comparing standard unsupervised training (C1) with label-based factor supervision in C2. 
Regime~D keeps the same correlated support and evaluation setting but replaces the monolithic latent with SP-VAE, allowing us to test whether explicit factor-addressable partitioning improves composition.
Each regime therefore provides a learned latent space in which the same conditional latent diffusion
architecture is trained. This allows the comparison to ask whether better latent factorisation leads
to more reliable compositional generation under a shared downstream setup.

\paragraph{Outcome interpretation.}
\begin{itemize}
    \item \textbf{Case 1:} If Regime~A succeeds, Regime~B fails, Regime~C1 partially recovers, Regime~C2 improves
    further, and Regime~D improves again, this supports the claim that supervision helps but
    factor-addressable partitioning helps more.
    \item \textbf{Case 2:} If Regime~A succeeds, Regime~B fails, Regime~C1 improves only modestly, and
    Regime~C2 already matches Regime~D, then label-based factor supervision may be sufficient for most of the
    desired factor separation in this setting.
    \item \textbf{Case 3:} If Regime~A succeeds, Regime~B fails, Regime~C1 and Regime~C2 remain limited, and
    Regime~D improves clearly, this supports the stronger interpretation that neither independence
    pressure nor supervision alone is sufficient under correlation without additional structural
    allocation.
    \item \textbf{Case 4:} If Regime~A already fails, the composition setup itself is not valid even in the control
    setting, and H2 is not supported.
\end{itemize}

\subsubsection{Regime A: Control}

Regime~A uses the standard dSprites dataset together with a baseline VAE, illustrated in Figure~\ref{fig:baseline_vae}. It tests whether
held-out recombination works when the data support is clean and factorised. The expected
outcome is that the model learns a representation that is sufficiently well aligned with
the underlying attributes for composition to recover both target factors.

The model in this regime is a vanilla variational autoencoder trained to encode dSprites
images into a continuous latent distribution and reconstruct them from that latent representation.
Its role in Phase~2 is to provide a standard generative model for testing how well a learned
latent space supports reconstruction and held-out attribute recombination.

The input data are the dSprites images from the standard split in Regime~A. The held-out
orientation--scale combination used for evaluation is removed from the training set, and the same
preprocessing is applied throughout Phase~2. Each example is represented as a single-channel
$64 \times 64$ sprite image.

\begin{figure}[H]
    \centering
    \includegraphics[width=\linewidth]{chapters/research_method/media/vae.drawio.pdf}
    \caption{Baseline VAE architecture used in Regimes~A and~B. An input dSprites image is encoded
    into the parameters of an approximate posterior, sampled via the reparameterization trick, and
    decoded into a reconstruction.}
    \label{fig:baseline_vae}
\end{figure}

Architecturally, the model follows the standard encoder--decoder VAE formulation. Given an input
image $x$, the encoder produces the parameters of an approximate posterior distribution,
$\mu(x)$ and $\log \sigma^2(x)$, from which a latent sample $z$ is drawn using the
reparameterization trick. The decoder then maps $z$ back to a reconstructed image $\hat{x}$ with
the same spatial dimensions as the input.

Its training objective is the negative ELBO, written as
\begin{align}
    \mathcal{L}_{\mathrm{VAE}}(x)
    &= \mathcal{L}_{\mathrm{rec}}(x, \hat{x}) + \mathcal{L}_{\mathrm{KL}}(x) \\
    &= -\mathbb{E}_{q_{\phi}(z \mid x)} \left[\log p_{\theta}(x \mid z)\right]
    + D_{\mathrm{KL}}\left(q_{\phi}(z \mid x)\,\|\,p(z)\right).
\end{align}
For binary dSprites images, the reconstruction term is the Bernoulli negative log-likelihood,
\begin{equation}
    \mathcal{L}_{\mathrm{rec}}(x, \hat{x})
    = -\sum_{i=1}^{64 \times 64}
    \left[x_i \log \hat{x}_i + (1 - x_i)\log(1 - \hat{x}_i)\right],
\end{equation}
and, assuming a diagonal Gaussian approximate posterior
$q_{\phi}(z \mid x) = \mathcal{N}\!\left(\mu(x), \mathrm{diag}(\sigma^2(x))\right)$ and a standard
normal prior $p(z) = \mathcal{N}(0, I)$, the KL term has the closed form
\begin{equation}
    \mathcal{L}_{\mathrm{KL}}(x)
    = \frac{1}{2}\sum_{j=1}^{d}
    \left(\mu_j(x)^2 + \sigma_j(x)^2 - \log \sigma_j(x)^2 - 1\right).
\end{equation}
The reconstruction term encourages $\hat{x}$ to match the observed image, while the KL term
regularises the approximate posterior toward the prior so that the model learns a smooth and
structured latent space.

In this way Regime~A establishes the positive control for Phase~2: if held-out recombination works
here, then the overall task is feasible before correlation is introduced. Regime~B keeps this same
architecture and ELBO objective fixed and changes only the data support.

\subsubsection{Regime B:  Correlated dSprites Setting}

Regime~B uses correlated dSprites with the same baseline VAE. It tests whether correlation in the
training distribution is sufficient to degrade the learned representation even when the image
domain, model, and evaluation procedure remain unchanged. The expected outcome is worse
held-out composition than in Regime~A, reflecting shortcut learning and latent entanglement
induced by the biased support.

Regime~B keeps the encoder-decoder architecture shown in
Figure~\ref{fig:baseline_vae}, the reparameterised latent code, and the ELBO objective from
Regime~A unchanged, but trains on the correlated dSprites split defined earlier. The held-out
orientation--scale combination is removed from training under the same protocol, so any
degradation relative to Regime~A can be attributed to the change in data support rather than to a
change in the model family or composition rule.

If performance degrades here, the failure setting is established and the next question becomes
whether benchmark disentanglement methods can recover any of the lost compositional performance
under the same correlated support.

\subsubsection{Regime C: Disentangling Correlated dSprites}
Regime~C is motivated by prior work showing that standard disentanglement objectives are brittle under correlated training data \citep{challeau2019challenging}, whereas weak supervision can mitigate this difficulty \citep{weakly_supervised_disentanglement}. 
\citet{traeuble2021disentangled} show that factorization-based inductive biases are insufficient when factors are correlated, and further show that weak supervision can help resolve the resulting latent entanglement. 
Regime~C asks whether composition under correlated support improves through stronger disentanglement objectives alone (C1), or through adding factor-label supervision while keeping the latent structure unchanged (C2).

Regime~C keeps the correlated dSprites dataset of Regime~B but replaces the baseline VAE with disentanglement models evaluated under two internal conditions:
\begin{itemize}
    \item \textbf{C1: Disentanglement.} Standard $\beta$-VAE \citep{Higgins2016betaVAELB} and
    FactorVAE~\citep{disentangling_by_factorising} test whether stronger independence
    objectives alone can recover compositional performance under correlation.
    \item \textbf{C2: Label-Based Factor Supervision.} We augment the C1 models with small
    auxiliary predictor networks (MLP heads) trained to recover the ground-truth scale (6-class)
    and orientation (40-class) labels from the encoder mean $\mu_\phi(x)$. The heads are
    training-only and are discarded at test
    time, leaving the generative model and composition interface identical to C1. This adds
    deterministic label-based guidance without altering the latent structure itself. To keep the
    comparison fair, we use the same labels and supervised loss terms as in Regime~D, isolating
    whether supervision alone improves composition.
\end{itemize}
We do not expect C1 to close the gap fully. 
$\beta$-VAE and FactorVAE encourage the latent space to separate factors, but this works best when those factors are not strongly correlated in the training data. 
Under correlation, that bias becomes less effective, so we expect only partial recovery in C1, consistent with \citet{traeuble2021disentangled}. 
C2 may improve further by adding factor-level supervision without changing the latent structure itself. 
If Regime~D still performs better, this would suggest that supervision alone is not enough, and that explicit factor-addressable partitioning provides an additional benefit.

At the representational and objective levels, these benchmark models retain a single latent code
$z \in \mathbb{R}^{d}$ and ask whether factor separation will emerge within that shared space
through stronger global regularisation:
\begin{equation}
\begin{aligned}
\text{$\beta$-VAE:}\qquad
& z \in \mathbb{R}^{d},
&& \mathcal{L}_{\beta\text{-VAE}}
= \mathcal{L}_{\mathrm{ELBO}} + \beta\,\mathcal{L}_{\mathrm{KL}}
\\[4pt]
\text{FactorVAE:}\qquad
& z \in \mathbb{R}^{d},
&& \mathcal{L}_{\mathrm{FactorVAE}}
= \mathcal{L}_{\mathrm{ELBO}} + \gamma\,\mathrm{TC}(q(z)).
\end{aligned}
\end{equation}
$\beta$-VAE
strengthens the KL bottleneck to encourage more factorised latent usage, while FactorVAE adds a
Total Correlation penalty that more directly discourages dependence across latent dimensions.

In \textbf{C1}, these models are evaluated in their standard unsupervised form. This condition
tests whether statistical independence objectives alone provide the right inductive bias for
composition under correlated support. Any improvement they achieve would indicate that generic
disentanglement pressure helps under bias; any remaining failure would suggest that improving independence alone does not guarantee that factors can be separately identified and recombined.

In \textbf{C2}, the same $\beta$-VAE and FactorVAE backbones are augmented with two
label-supervised auxiliary predictors $g_{\mathrm{scale}}$ and $g_{\mathrm{orient}}$, each a small
two-layer MLP, attached to the encoder mean $\mu_\phi(x)$ rather than to the stochastic latent
sample~$z$. Using $\mu_\phi(x)$ makes the supervised signal deterministic: the auxiliary
gradients do not flow through the reparameterisation noise $\varepsilon$, which reduces variance
in the label-supervision signal and is consistent with standard disentanglement evaluation
practice that assesses factor information through $\mu$ rather than
$z$~\citep{eastwood2018framework, challeau2019challenging}. The two predictors produce logits
\begin{equation}
\hat{y}_s = g_{\mathrm{scale}}\!\left(\mu_\phi(x)\right) \in \mathbb{R}^6,
\qquad
\hat{y}_o = g_{\mathrm{orient}}\!\left(\mu_\phi(x)\right) \in \mathbb{R}^{40},
\end{equation}
and are trained with cross-entropy against the dSprites dataset labels
$y_s \in \{0,\ldots,5\}$ and $y_o \in \{0,\ldots,39\}$ for scale and orientation respectively.
The full C2 objective is
\begin{equation}
\label{eq:c2_objective}
\mathcal{L}_{\mathrm{C2}}
= \mathcal{L}_{\mathrm{backbone}}
+ \lambda_s\,\mathrm{CE}(\hat{y}_s,\, y_s)
+ \lambda_o\,\mathrm{CE}(\hat{y}_o,\, y_o),
\end{equation}
where $\mathcal{L}_{\mathrm{backbone}}$ is either the $\beta$-VAE or FactorVAE objective defined
above, and $\lambda_s$, $\lambda_o$ are scalar weights matched to those used in Regime~D.
Figure~\ref{fig:c2_architecture} then illustrates the complete C2 forward pass.

\begin{figure}[H]
    \centering
    \includegraphics[width=\linewidth]{chapters/research_method/media/c2_architecture.pdf}
    \caption{C2 architecture. The $\beta$-VAE or FactorVAE backbone is unchanged. Two auxiliary
    MLP heads ($g_{\mathrm{scale}}$, $g_{\mathrm{orient}}$) receive the encoder mean
    $\mu_\phi(x)$ and are trained with cross-entropy against the 6-class scale and 40-class
    orientation labels introduced in Figure~\ref{fig:c2_factor_targets}. Both heads are discarded
    at test time; composition proceeds in the unchanged shared latent space~$z$.}
    \label{fig:c2_architecture}
\end{figure}

Both auxiliary heads are training-only and are discarded at test time: the composition rule
operates in the same shared latent space $z$ as in Regimes~A and~C1, without modification.
This follows the established pattern of training-only auxiliary branches in representation
learning, where removing the auxiliary head at inference is precisely the mechanism by which it
shapes the main encoder representation without coupling the downstream interface to the
supervision signal~\citep{chen2020simple}.

C2 is architecturally distinct from the M2 semi-supervised VAE~\citep{kingma2014semi}, which
conditions the decoder on the label $y$ and marginalises over it during inference, thereby
changing the generative model structure. In C2, the decoder sees only $z$ and labels never
enter the generative path. Labels are used solely as a training-time signal to shape the
encoder's representation of $\mu_\phi(x)$, keeping the generative model and the composition
interface identical to Regimes~A and~B.

These adaptations use the same factor labels, prediction losses, and predictor-loss weights
later used in SP-VAE, following the broader logic of weakly supervised disentanglement
comparisons in which limited factor supervision is introduced explicitly rather than left
implicit in the latent structure~\citep{locatello2020weakly}. Regime~C2 therefore serves as a
supervision-only control: if it matches SP-VAE, then additional structure may not be necessary;
if SP-VAE improves beyond C2, the remaining gain is more plausibly attributable to the explicit
latent partitioning introduced in Regime~D rather than to the shared supervision signal.

\subsubsection{Regime D: Partitioned Representation}
\label{subsec:p2_spvae}

Regimes~A--C all operate in a single shared latent space, $z \in \mathbb{R}^d$. When downstream
composition requires orientation and scale to act as separable conditioning signals, that shared
representation does not guarantee a clean separation: if these factors overlap in the same latent
dimensions, the model can still represent them jointly rather than compositionally. SP-VAE
addresses this by assigning preserved context, orientation, and scale to separate latent parts,
making that separation more explicit.

Regime~D keeps the correlated dSprites support, the held-out combinations, and the same downstream
evaluation setting used in Regimes~B and~C, but changes the representation itself.
This makes it the clearest Phase~2 test of H2: when supervision is matched to Regime~C2 and the downstream evaluation architecture and inference
procedure are held fixed, any additional gain is most directly attributable to how the
latent variables are organised rather than to a different inference mechanism.

SP-VAE is inspired by SepVAE's common-versus-salient partitioning~\citep{louiset2024sepvae}, but
is adapted here to Phase~2's single-dataset, multi-factor setting. SepVAE cannot be used unchanged
because its contrastive background/target split and associated objectives rely on a dataset
asymmetry that correlated dSprites does not provide. SP-VAE therefore keeps the C2 supervision
setting but replaces the shared latent with a partitioned code
\(
z=(z_c,z_{d1},z_{d2}),
\)
where \(z_c\) stores shared context, \(z_{d1}\) stores orientation, and \(z_{d2}\) stores scale.
Figure~\ref{fig:spvae_architecture} summarises this architecture.


\begin{figure}[H]
    \centering
    \includegraphics[width=\linewidth]{chapters/research_method/media/spvae_architecture.pdf}
    \caption{SP-VAE architecture for Regime~D. A shared encoder produces three posterior heads for shared
context, orientation, and scale. Their samples are concatenated and decoded to reconstruct
\(\hat{x}\). Training-only classifiers supervise the intended branch and probe the opposite branch
through GRL.}
    \label{fig:spvae_architecture}
\end{figure}

Unlike $\beta$-VAE, FactorVAE, and Regime~C2, which all retain a single shared latent and try to
organize factors within it, SP-VAE separates shared context, orientation, and scale into different
latent parts. This matters especially under correlated support: when orientation and scale co-occur
systematically in training, a shared latent can still exploit that shortcut and store them
together, even under stronger regularization or label supervision. SP-VAE addresses this in three
steps: first, it assigns each role to a separate latent part; second, it uses training-time
classifiers to align each factor-specific part with its intended attribute; and third, it applies
adversarial wrong-head supervision to discourage redundant leakage across those parts. The result is
a representation designed not just to be cleaner in the abstract, but to support the later
substitution operation directly.

Figure~\ref{fig:spvae_architecture} shows the resulting SP-VAE architecture. Given an input image
$x$, the encoder produces three posterior parameter heads,
\((\mu_c, \log \sigma_c^2)\), \((\mu_{d1}, \log \sigma_{d1}^2)\), and
\((\mu_{d2}, \log \sigma_{d2}^2)\), corresponding to shared context, orientation, and scale.
Reparameterized samples $z_c$, $z_{d1}$, and $z_{d2}$ are then concatenated into \(z\) and passed
to the decoder to reconstruct \(\hat{x}\).

During training, supervision is applied directly to the encoder means of the two factor latents.
The orientation classifier takes \(\mu_{d1}\) as input and produces 40 orientation logits, while
the scale classifier takes \(\mu_{d2}\) as input and produces 6 scale logits. These logits are
compared with the ground-truth labels using cross-entropy, which pushes \(\mu_{d1}\) to encode
orientation and \(\mu_{d2}\) to encode scale. We use the posterior means rather than sampled
latents so that this supervised signal remains stable and is not perturbed by reparameterization
noise, matching the supervision setup already used in Regime~C2.

Correct-head supervision alone does not guarantee exclusivity: even if
\(g_{\mathrm{orient}}(\mu_{d1})\) predicts orientation well and
\(g_{\mathrm{scale}}(\mu_{d2})\) predicts scale well, orientation may still be encoded in
\(\mu_{d2}\) and scale may still be encoded in \(\mu_{d1}\) if that helps reconstruction. In other
words, the positive cross-entropy terms push the right factor into the right latent part, but they
do not push the wrong factor out of the wrong one. An additional mechanism is therefore needed to
actively suppress this cross-factor leakage.

To provide that suppression, we apply each classifier to the opposite latent mean through a
gradient reversal layer, giving the wrong-head branches
\(g_{\mathrm{orient}}(\mathrm{GRL}(\mu_{d2}))\) and
\(g_{\mathrm{scale}}(\mathrm{GRL}(\mu_{d1}))\). The GRL is identity in the forward pass, so the
classifier receives the real latent activation and produces an ordinary cross-entropy loss, but in
the backward pass it reverses the gradient sent to the encoder~\citep{ganin2015unsupervised,
ganin2016domain}. Without this reversal, the wrong-head loss would encourage the encoder to make the
non-target factor easier to predict, that is, the unreversed gradient would push the encoder to
encode that factor more strongly in the wrong latent part. 
With GRL, the classifier still tries to predict the non-target factor from the wrong latent part.
The encoder, however, receives the reversed gradient and is therefore trained to make that
prediction harder. In this way, the classifier and encoder are optimized in opposite directions: the
classifier becomes better at detecting any leakage that remains, while the encoder becomes better at
removing it. This is the same basic idea used in domain-adversarial training, where gradient
reversal is used to remove domain information from features; here, the unwanted signal is not domain
identity, but the wrong factor~\citep{attribute_disentanglement}.

The final SP-VAE objective consists of four parts: the ELBO, correct-head supervision, wrong-head suppression, and residual separation penalties:
\begin{equation}
\begin{aligned}
\mathcal{L}_{\mathrm{SP\mbox{-}VAE}}
&= \mathcal{L}_{\mathrm{rec}}(x,\hat{x})
+ \beta_c \, D_{\mathrm{KL}}\!\left(q_\phi(z_c \mid x)\,\|\,p(z_c)\right)
+ \beta_{d1} \, D_{\mathrm{KL}}\!\left(q_\phi(z_{d1} \mid x)\,\|\,p(z_{d1})\right) \\
&\quad
+ \beta_{d2} \, D_{\mathrm{KL}}\!\left(q_\phi(z_{d2} \mid x)\,\|\,p(z_{d2})\right)
+ \lambda^{+}_{o}\,\mathrm{CE}(g_{\mathrm{orient}}(\mu_{d1}), y_o)
+ \lambda^{+}_{s}\,\mathrm{CE}(g_{\mathrm{scale}}(\mu_{d2}), y_s) \\
&\quad
+ \lambda^{-}_{o}\,\mathrm{CE}(g_{\mathrm{orient}}(\mathrm{GRL}(\mu_{d2})), y_o)
+ \lambda^{-}_{s}\,\mathrm{CE}(g_{\mathrm{scale}}(\mathrm{GRL}(\mu_{d1})), y_s) \\
&\quad
+ \lambda_{\mathrm{mi}} \hat{I}(z_{d1}; z_{d2})
+ \lambda_{\mathrm{orth}} \lVert \mathrm{Cov}(z_{d1}, z_{d2}) \rVert_F^2 .
\end{aligned}
\end{equation}

where:
\begin{itemize}
    \item \(\mathcal{L}_{\mathrm{ELBO}}\) is the standard variational objective, combining
    reconstruction and KL regularization. It keeps SP-VAE a generative model and preserves the
    sampling and reconstruction behaviour expected of a VAE.

    \item \(\lambda^{+}_{o}\,\mathrm{CE}(g_{\mathrm{orient}}(\mu_{d1}), y_o)\) and
    \(\lambda^{+}_{s}\,\mathrm{CE}(g_{\mathrm{scale}}(\mu_{d2}), y_s)\) are the correct-head
    supervision terms. Following weakly supervised disentanglement and attribute-specific
    supervision, they align each factor latent with its intended label, pushing orientation into
    \(z_{d1}\) and scale into \(z_{d2}\)~\citep{weakly_supervised_disentanglement,
    attribute_disentanglement}.

    \item \(\lambda^{-}_{o}\,\mathrm{CE}(g_{\mathrm{orient}}(\mathrm{GRL}(\mu_{d2})), y_o)\) and
    \(\lambda^{-}_{s}\,\mathrm{CE}(g_{\mathrm{scale}}(\mathrm{GRL}(\mu_{d1})), y_s)\) are the
    wrong-head suppression terms. Through gradient reversal, they encourage the encoder to remove
    orientation from \(z_{d2}\) and scale from \(z_{d1}\), following domain-adversarial training and
    GRL-based attribute disentanglement~\citep{ganin2015unsupervised, ganin2016domain,
    attribute_disentanglement}.

    \item \(\lambda_{\mathrm{mi}} \hat{I}(z_{d1}; z_{d2})\) penalizes residual statistical
    dependence between the two factor latents. This follows the broader disentanglement literature on
    reducing dependence between latent subspaces, and is included here because classifier-based terms
    may still leave weaker correlations behind~\citep{disentangling_by_factorising,
    funke2022disentanglement}.

    \item \(\lambda_{\mathrm{orth}} \lVert \mathrm{Cov}(z_{d1}, z_{d2}) \rVert_F^2\) penalizes
    cross-covariance between the factor latents, reducing geometric overlap between their subspaces.
    This complements the mutual-information term by discouraging linear overlap even when residual
    dependence is too weak to be detected reliably by the classifiers~\citep{disentengaled_representation_learning}.
\end{itemize}

This design is meant to produce a latent space in which orientation and scale are cleanly separated
before they are passed to the shared downstream conditional diffusion evaluator. Direct slot
substitution therefore serves as an architecture-specific sanity check on SP-VAE's interpretability,
while the main cross-regime evaluation uses the common conditional diffusion setup described below.
We therefore design the representation to match the operation
we want to perform on it. The expected outcome is stronger held-out composition than in Regime~C2,
which would support H2 by showing that giving preserved context and each target factor its own
cleanly organized latent part makes recombination easier.

\subsection{Evaluation}
\label{subsec:p2_eval}

Phase~2 evaluates each regime as a downstream stress test of representational usefulness. First, we
measure how well the learned latent code isolates the known dSprites factors. Second, we train the
same conditional latent diffusion architecture in each regime's learned latent space and test
whether that shared downstream setup can recover held-out orientation-scale combinations. The same
external attribute classifiers are then used to score every generated image.

\paragraph{Attribute Classifiers.}
Two attribute classifiers are trained on data splits disjoint from VAE training:
\begin{itemize}
    \item $f_{\mathrm{or}}$: predicts the orientation value of an input image.
    \item $f_{\mathrm{sc}}$: predicts the scale value of an input image.
\end{itemize}
These classifiers are external evaluators, not training-time heads. Given a generated image,
$f_{\mathrm{or}}$ predicts its orientation and $f_{\mathrm{sc}}$ predicts its scale. Their
predictions are used to determine whether the generated sample has recovered the intended target
factors.

Each composed image is evaluated with the orientation and scale classifiers. The composition success
rate is the proportion of held-out source pairs for which the generated image is classified as
recovering both the target orientation and the target scale.

% \paragraph{Disentanglement Metric.}
% H2 claims that composition success is limited by representation quality, so we also report a
% standard disentanglement diagnostic alongside the output-level composition results. We use
% DCI~\citep{eastwood2018framework}, computed on the deterministic encoder representation:
% $\mu_\phi(x)$ for Regimes~A--C and $[\mu_c,\mu_{d1},\mu_{d2}]$ for Regime~D. Operationally, DCI fits
% predictors from the latent code to the known dSprites factors and summarises how selectively that
% information is organised in the representation.

% \begin{itemize}
%     \item \textbf{Disentanglement}: whether each latent dimension is associated mainly with one
%     factor rather than several.
%     \item \textbf{Completeness}: whether each factor is concentrated in a small number of latent
%     dimensions rather than spread broadly across the code.
%     \item \textbf{Informativeness}: whether the known factors are recoverable from the latent code
%     at all.
% \end{itemize}

% Higher disentanglement and completeness therefore indicate cleaner factor separation, while
% informativeness confirms that the code still preserves usable information. These scores are reported
% for all regimes as explanatory evidence for why composition succeeds or fails under correlated
% support

\paragraph{Held-out composition test.}
For each regime $r$, the learned encoder-decoder pair defines a latent space in which images can be
encoded and decoded. We then train the same conditional latent diffusion architecture in that latent
space, using orientation and scale labels as conditioning variables. At inference time, composition
is performed with the same score-composition rule in every regime, combining the unconditional,
orientation-conditioned, and scale-conditioned denoising estimates following the standard
compositional diffusion literature~\citep{liu2022compositional, du2024compositional}. The resulting
latent sample is decoded with the regime's decoder to produce the composed image $x_{\mathrm{comp}}$.
\begin{equation}
\hat{\epsilon}^{\,\mathrm{comp}}_r(z_t,t)
=
\hat{\epsilon}_r(z_t,t,\emptyset)
+
w_{\mathrm{or}}\bigl(\hat{\epsilon}_r(z_t,t,y_{\mathrm{or}})-\hat{\epsilon}_r(z_t,t,\emptyset)\bigr)
+
w_{\mathrm{sc}}\bigl(\hat{\epsilon}_r(z_t,t,y_{\mathrm{sc}})-\hat{\epsilon}_r(z_t,t,\emptyset)\bigr).
\end{equation}
Each composed image is evaluated with the orientation and scale classifiers. The composition success
rate is the proportion of held-out source pairs for which the generated image is classified as
recovering both the target orientation and the target scale.

% \paragraph{Stress Test Limitation}
% Because each regime defines a different latent geometry, each regime also requires a separately
% trained conditional diffusion model. The comparison is therefore controlled at the level of
% architecture, conditioning scheme, inference procedure, and output scoring, but it does not isolate
% representation alone in the strongest causal sense. Phase~2 should therefore be read as a shared
% downstream stress test of whether better latent factorisation supports compositional diffusion more
% reliably.

\paragraph{Trajectory-level gap closure.}
We also compare the denoising trajectory under composed
conditioning with the trajectory obtained from a single joint condition on the same target pair. The
terminal-gap diagnostic $d_T$ is then used as a secondary indicator of whether the
composability gap narrows under improved representation learning: smaller terminal gaps indicate
that the composed trajectory tracks the jointly conditioned trajectory more closely.

\subsection{Summary}
\label{subsec:p2_bridge}

Phase~2 addresses the gap identified in Phase~1 by testing H2 in a controlled synthetic setting.
Regime~A provides the control, Regime~B creates the failure setting, Regimes~C1 and~C2 test
whether stronger disentanglement or supervision recover performance, and Regime~D tests whether
explicit factor structure helps further. If performance drops from Regime~A to Regime~B and then
improves through Regimes~C1, C2, and especially~D, Phase~2 supports H2 by showing that better
representations can improve downstream compositional generation and help reduce the gap in this
controlled setting. Conversely, if Regimes~C1, C2, and~D do not improve clearly over Regime~B, then Phase~2 does not support H2: in this setting, learning a more factorised latent space is not sufficient to improve composition, and any further gains may require changes to the composition rule or other inference-time mechanisms.

Phase~3 then asks whether the same idea remains useful for chest X-rays, where
conditional independence is approximate and success must be judged by radiological plausibility.


\begin{table}[H]
    \centering
    \small
    \caption{Summary of the Phase~2 representation regimes. Regime~C is split into C1 and C2 so
    that standard disentanglement objectives, label-based factor supervision, and structured latent
    partitioning can be compared under the same correlated dSprites support and held-out
    composition test.}
    \label{tab:p2_representation_regimes}
    \begin{tabularx}{\linewidth}{|l|p{4.2cm}|X|X|}
        \hline
        Regime & Model and training condition & Question & Expected outcome \\
        \hline
        A & Baseline VAE trained on standard dSprites with the ELBO only. & Can held-out
        recombination succeed under clean, factorised support? & High composition success in the
        control setting. \\
        \hline
        B & Baseline VAE trained on correlated dSprites with the ELBO only. & Does correlation in
        the training support induce entanglement and degrade held-out composition? & Lower
        composition success than Regime~A. \\
        \hline
        C1 & Standard $\beta$-VAE and FactorVAE trained on correlated dSprites with their native
        disentanglement objectives and no auxiliary predictors. & Are standard disentanglement
        objectives alone sufficient to recover performance under correlation? & Partial recovery
        relative to Regime~B if independence pressure helps. \\
        \hline
        C2 & $\beta$-VAE and FactorVAE trained on correlated dSprites with label-based auxiliary
        predictors for scale and orientation in addition to their native objectives. & Does label-based factor supervision
        improve composition without changing the monolithic latent structure? & Additional recovery
        relative to C1 if supervision alone helps. \\
        \hline
        D & SP-VAE trained on correlated dSprites with the same auxiliary factor supervision, explicit
        latent partitioning, and the additional separation constraints introduced below. & Does
        factor-addressable latent allocation add beyond label-based factor supervision alone? & Stronger
        recovery than Regime~C if structured partitioning matters. \\
        \hline
    \end{tabularx}
\end{table}
